% Chapter 14: What Does Not Imply What
% This chapter is witness-centered: each separation lives or dies by
% the construction that demonstrates it.
% Content: ALL 9 does_not_imply edges, ALL 4 strictly_stronger edges.

\chapter{What Does Not Imply What}
\label{ch:separations}

Most textbooks in learning theory treat separation results as scattered
remarks: brief asides after a characterization theorem, noting that
some plausible implication fails. This chapter reverses that convention.
Here, the separations are first-class citizens: each one receives a
formal statement, a witness construction, and an analysis of what
structural feature the witness exploits.

A separation result has two components. The \emph{statement} asserts
that some implication $A \Rightarrow B$ does not hold. The
\emph{witness} is a concrete mathematical object (a concept class,
a dimension pair, a computational reduction) that demonstrates the
failure. The witness is the mathematics; the statement is merely its
summary. Throughout this chapter, we privilege the construction over the
claim.

We organize the chapter's 13 edges into two groups: 9
\rel{does\_not\_imply} edges, where the non-implication is the content,
and 4 \rel{strictly\_stronger} edges, where an implication does hold but
is provably non-reversible. Together they form the \emph{separation
lattice} of formal learning theory.

% ================================================================
% SECTION 1: THE SEPARATION LATTICE
% ================================================================

\section{The Separation Lattice}
\label{sec:sep-lattice}

\Cref{fig:separation-lattice} displays all 13 edges as a single
diagram. Dashed red arrows denote \rel{does\_not\_imply} edges: the
source does \emph{not} entail the target, and the label names the
witness. Solid blue arrows denote \rel{strictly\_stronger} edges: the
source strictly contains the target as a special case, with the witness
demonstrating the gap.

The first observation is structural: the lattice is \emph{sparse}.
Thirteen edges connect concepts drawn from six paradigms and a dozen
complexity measures. Most paradigm pairs are simply incomparable: they
neither imply nor contradict each other, because they operate on
different mathematical objects. The sparsity is itself informative:
learning theory is not a linear hierarchy from weak to strong, but a
partially ordered collection of largely independent formalisms.

\begin{figure}[ht]
\centering
\begin{tikzpicture}[
  node distance=1.8cm and 2.4cm,
  concept/.style={
    rectangle, rounded corners=3pt, draw, thick,
    fill=layergray, font=\small\sffamily,
    minimum width=2.2cm, minimum height=0.7cm,
    align=center
  },
  dni/.style={-{Stealth[length=5pt]}, dashed, thmred, thick},
  ss/.style={-{Stealth[length=5pt]}, defblue, thick},
  witness/.style={font=\tiny\itshape, text=notegray, align=center}
]

% --- Nodes ---
\node[concept] (ul) {Universal\\Learning};
\node[concept, below left=2.0cm and 1.2cm of ul] (ol) {Online\\Learning};
\node[concept, below right=2.0cm and 1.2cm of ul] (pac) {PAC\\Learning};
\node[concept, right=2.8cm of pac] (exact) {Exact\\Learning};
\node[concept, below=1.8cm of ol] (mb) {Mistake\\Bounded};
\node[concept, below=1.8cm of pac] (ex) {Ex-\\Learning};
\node[concept, below=1.8cm of ex] (fin) {FIN-\\Learning};

\node[concept, left=2.5cm of ol] (vc) {VC\\Dimension};
\node[concept, below=1.8cm of vc] (sq) {SQ\\Dimension};
\node[concept, below left=1.8cm and 0.0cm of vc] (comp) {Comput.\\Hardness};

\node[concept, right=2.8cm of exact] (nd) {Natarajan\\Dimension};
\node[concept, above=1.8cm of nd] (ds) {DS\\Dimension};

\node[concept, below=1.8cm of mb] (ft) {Fundamental\\Theorem};
\node[concept, right=2.8cm of ft] (stab) {Algorithmic\\Stability};

\node[concept, below left=1.4cm and 0.5cm of pac] (pi) {Proper/Improper\\Separation};
\node[concept, below right=1.8cm and 0.0cm of exact] (uc) {Unlabeled\\Compression};
\node[concept, below=1.5cm of uc] (cs) {Compression\\Scheme};

% --- does_not_imply edges (dashed red) ---
\draw[dni] (pac) -- node[witness, above, sloped] {thresholds} (mb);
\draw[dni] (nd) -- node[witness, above, sloped] {$d_N{=}1$, not learnable} (pac);
\draw[dni] (pac) -- node[witness, above, sloped] {approx.\ vs.\ exact} (exact);
\draw[dni] (vc) -- node[witness, left, pos=0.4] {crypto} (comp);
\draw[dni] (vc) -- node[witness, left] {parities} (sq);
\draw[dni] (ex) -- node[witness, above, sloped] {$\VC = \infty$} (pac);
\draw[dni] (pi) -- node[witness, below, sloped] {3-DNF} (pac);
\draw[dni] (uc) -- node[witness, right, pos=0.4] {P\'alv\"olgyi--Tardos} (cs);
\draw[dni] (ft) -- node[witness, above, sloped] {beyond binary} (stab);

% --- strictly_stronger edges (solid blue) ---
\draw[ss] (ul) -- node[witness, above left, sloped] {removes uniformity} (pac);
\draw[ss] (ol) -- node[witness, above, sloped] {$\Ldim \Rightarrow \VC$} (pac);
\draw[ss] (ex) -- node[witness, right] {unbounded mind changes} (fin);
\draw[ss] (ds) -- node[witness, right] {pseudo-manifold} (nd);

\end{tikzpicture}

\caption{The separation lattice. Dashed red:
\rel{does\_not\_imply} (9~edges). Solid blue:
\rel{strictly\_stronger} (4~edges). Each label names the witness.}
\label{fig:separation-lattice}
\end{figure}


% ================================================================
% SECTION 2: SEPARATIONS BETWEEN PARADIGMS
% ================================================================

\section{Separations Between Paradigms}
\label{sec:separations}

We now present the 9 \rel{does\_not\_imply} edges in the graph, each
with its witness construction. The order moves from the most elementary
witness (a one-parameter family on $\R$) to the most conceptually
involved (the breakdown of the fundamental theorem beyond binary
classification).

% --- Separation 1 ---

\begin{separation}
\textbf{PAC learning $\not\Rightarrow$ mistake-bounded learning.}\formal{FLT_Proofs.Theorem.Separation.pac_not_implies_online}

\medskip\noindent
\emph{Witness.}
Let $\mathcal{H} = \{h_\theta : \theta \in \R\}$ be the class of
thresholds on $\R$, where $h_\theta(x) = \ind[x \geq \theta]$.

This class has $\VC(\mathcal{H}) = 1$: a single point $x$ is shattered
(choose $\theta < x$ or $\theta > x$), but no two points can be
simultaneously shattered.\formal{FLT_Proofs.Theorem.Separation.vcdim_threshold_finite} By the fundamental theorem of statistical
learning, $\mathcal{H}$ is PAC learnable with sample complexity
$O(1/\varepsilon)$.

However, $\Ldim(\mathcal{H}) = \infty$.\formal{FLT_Proofs.Theorem.Separation.ldim_threshold_top} An adaptive adversary can force
arbitrarily many mistakes by binary search: maintain an interval
$(a, b)$ of uncertainty, present the midpoint, and whichever label the
learner predicts, place the true threshold on the opposite side. Each
round forces a mistake and halves the interval, but the reals admit
no finite termination.

\medskip\noindent
\emph{What the witness exploits.}
The gap between VC and Littlestone dimension: batch learnability
(distribution-free, i.i.d.\ samples) does not imply sequential
learnability (adversarial, adaptive instances). The reals are
``too dense'' for any online strategy to pin down the threshold,
but a random sample reveals it with high probability.

\hfill\cite{Littlestone1988}
\end{separation}

% --- Separation 2 ---

\begin{separation}
\textbf{Ex-learning $\not\Rightarrow$ PAC learning.}\formal{FLT_Proofs.Theorem.Separation.ex_not_implies_pac}

\medskip\noindent
\emph{Witness.}
Let $\mathcal{C} = \{ S \subseteq \N : S \text{ is finite} \}$, the
class of all finite subsets of $\N$. A Gold-style learner can identify
any finite $S$ in the limit from positive data: enumerate elements as
they appear, and at each stage conjecture the set of elements seen so
far. After the last new element arrives, the conjecture stabilizes on
$S$.

However, $\VC(\mathcal{C}) = \infty$: for any $n$ points
$\{x_1, \ldots, x_n\} \subset \N$, every subset is itself a finite set
in $\mathcal{C}$, so the class shatters every finite set. The
fundamental theorem then implies $\mathcal{C}$ is not PAC learnable.

\medskip\noindent
\emph{What the witness exploits.}
The data models are incompatible. Gold-style identification receives an
infinite enumeration of the target's elements and succeeds in the limit;
PAC learning receives a finite i.i.d.\ sample and must generalize
immediately. A class can be identifiable from exhaustive enumeration yet
have no finite VC dimension.

\hfill\cite{Gold1967}
\end{separation}

% --- Separation 3 ---

\begin{separation}
\textbf{PAC learning $\not\Rightarrow$ exact learning.}

\medskip\noindent
\emph{Witness.}
Exact learning in the sense of Angluin~\cite{Angluin1988} identifies the
target \emph{exactly}, with zero error, from membership and equivalence
queries.  It departs from PAC learning along \emph{two} axes at once: the
success criterion (zero error versus $\varepsilon$-approximate) and the
data model (active queries versus a passive i.i.d.\ sample).  A
PAC-learnable class need not be exactly learnable on either count.

The criterion axis alone already bites computationally. Consider a class
$\mathcal{C}$ that is PAC learnable (finite VC dimension) but for which
finding a hypothesis \emph{exactly} consistent in the proper representation
is NP-hard (intersections of halfspaces in $\R^d$ are the standard
example~\cite{PittValiant1988}). An improper PAC learner attains low error
through a surrogate class, yet exact identification of the target is
intractable.

\medskip\noindent
\emph{What the witness exploits.}
Two mismatches stacked: approximate-versus-exact in the criterion, and
i.i.d.-sample-versus-query in the data model. PAC tolerates
$\varepsilon$~error and reads a passive sample; exact learning tolerates
none and may interrogate the target. Neither model subsumes the other.

\hfill\cite{Angluin1988, PittValiant1988}
\end{separation}

% --- Separation 4 ---

\begin{separation}
\textbf{Finite VC dimension $\not\Rightarrow$ computational
tractability.}

\medskip\noindent
\emph{Witness.}
Kearns and Valiant~\cite{KearnsValiant1994} constructed concept classes
based on polynomial-size Boolean circuits whose VC dimension is
polynomial in the circuit size, yet for which PAC learning is
computationally intractable under the assumption that the
\emph{decisional Diffie--Hellman} (or related cryptographic) assumption
holds.

The construction proceeds as follows. Fix a one-way function family.
Define $\mathcal{C}$ to be the class of functions computed by circuits
of size $s$. This class has $\VC(\mathcal{C}) \leq O(s \log s)$, so it
is information-theoretically PAC learnable with $\mathrm{poly}(s)$
samples. But any polynomial-time PAC learner for $\mathcal{C}$ could be
used to invert the one-way function, contradicting the cryptographic
assumption.

\medskip\noindent
\emph{What the witness exploits.}
The fundamental theorem is \emph{information-theoretic}: it
characterizes learnability in terms of sample complexity, not
computational complexity. Cryptographic hardness lives in a different
layer entirely. Finite VC dimension guarantees that enough data suffices;
it says nothing about whether a polynomial-time algorithm can find a
good hypothesis from that data.

\hfill\cite{KearnsValiant1994}
\end{separation}

% --- Separation 5 ---

\begin{separation}
\textbf{Finite VC dimension $\not\Rightarrow$ low SQ dimension.}

\medskip\noindent
\emph{Witness.}
The class of parities over $\{0,1\}^n$. Each parity function $\chi_S$,
for $S \subseteq [n]$, maps $x \mapsto \bigoplus_{i \in S} x_i$. The
parity class has $\VC = n$ (it shatters the standard basis). But any
statistical query algorithm requires queries of tolerance
$\tau = O(2^{-n/2})$ or uses $2^{\Omega(n)}$ queries of polynomial
tolerance, because distinct parities are pairwise uncorrelated under the
uniform distribution.

Formally, $\SQdim(\text{Parities}_n) = 2^n$, which is exponential in
$\VC = n$.

\medskip\noindent
\emph{What the witness exploits.}
VC dimension measures combinatorial shattering capacity
(distribution-free). SQ dimension measures pairwise correlation
structure under a specific distribution. Parities are maximally
uncorrelated under the uniform distribution, forcing any SQ learner to
make exponentially many queries, even though the class is small in the
VC sense.

\hfill\cite{BlumFurstJacksonKearnsMansourRudich1994}
\end{separation}

% --- Separation 6 ---

\begin{separation}
\textbf{Natarajan dimension $\not\Rightarrow$ multiclass PAC
learnability.}

\medskip\noindent
\emph{Witness.}
Brukhim et al.~\cite{BrukhimCarmonDinurMoranYehudayoff2022}
constructed a concept class $\mathcal{C}$ with label space
$Y$ of infinite cardinality, Natarajan dimension $\Ndim(\mathcal{C})=1$,
yet $\mathcal{C}$ is not PAC learnable.

The Natarajan dimension, which
generalizes VC dimension to multiclass settings by requiring two
distinct labelings of shattered points, is too coarse when $|Y|$ is
infinite. The construction builds a class in which every pair of points
can be 2-colored in only one way (so $\Ndim = 1$), but the global
combinatorial structure is rich enough to prevent uniform convergence.

\medskip\noindent
\emph{What the witness exploits.}
The Natarajan dimension captures pairwise shattering. When the label
space is infinite, pairwise structure does not determine global
learnability. The DS dimension, which accounts for higher-order
combinatorial structure via oriented hypergraphs, is the correct
characterization.

\hfill\cite{BrukhimCarmonDinurMoranYehudayoff2022}
\end{separation}

% --- Separation 7 ---

\begin{separation}
\textbf{Proper learnability $\not\Rightarrow$ efficient proper PAC
learning.}

\medskip\noindent
\emph{Witness.}
The class of 3-term DNF formulas over $\{0,1\}^n$. Pitt and
Valiant~\cite{PittValiant1988} showed that \emph{proper} PAC
learning\formal{FLT_Proofs.Basic.IsProper} (where the hypothesis must
itself be a 3-term DNF) is NP-hard, assuming
$\mathrm{RP} \neq \mathrm{NP}$. Yet the same class
is efficiently PAC learnable \emph{improperly}: every 3-term DNF can be
represented as a 3-CNF, and 3-CNFs are efficiently PAC learnable.

\medskip\noindent
\emph{What the witness exploits.}
The gap between proper and improper learning is purely
computational: the information-theoretic sample complexity is
identical, but the representation constraint of proper learning
introduces NP-hardness. The concept class is simple enough to learn
with the right representation, but finding a hypothesis in the
\emph{same} representation class is intractable.

\hfill\cite{PittValiant1988}
\end{separation}

% --- Separation 8 ---

\begin{separation}
\textbf{Labeled compression $\not\Rightarrow$ unlabeled compression.}

\medskip\noindent
\emph{Witness.}
P\'alv\"olgyi and Tardos~\cite{PalvolgyiTardos2020} exhibited a
concept class with $\VC = 2$ that admits a labeled sample compression
scheme of size~2 but does \emph{not} admit an unlabeled compression
scheme of size~2.

In an unlabeled compression scheme, the compression function stores only
the \emph{points} (not their labels) from the training sample; the
reconstruction function must infer both the hypothesis and the labels
from the point set alone. The witness class is constructed so that the
label information is essential for reconstruction: two subsets of the
same size can correspond to different concepts, and only the labels
disambiguate them.

\medskip\noindent
\emph{What the witness exploits.}
The label information carries entropy that cannot always be recovered
from the geometry of the point set alone. Unlabeled compression demands
that the point configuration determines the concept, which is a strictly
stronger requirement than labeled compression.

\hfill\cite{PalvolgyiTardos2020}
\end{separation}

% --- Separation 9 ---

\begin{separation}
\textbf{The fundamental theorem $\not\Rightarrow$ stability
characterization.}

\medskip\noindent
\emph{Witness.}
The Fundamental Theorem of Statistical Learning (\Cref{ch:pac}) ties
together VC dimension, uniform convergence, and PAC learnability into a
single equivalence.  Nine conditions, all the same condition.  That
theorem is the crown jewel of binary classification.

It does not survive contact with real-valued loss.

Shalev-Shwartz et al.~\cite{ShalevShwartzShamir2010} construct a
learning problem (arbitrary loss, multi-valued output) that is
learnable, yet for which uniform convergence fails outright.  The
nine-way equivalence collapses.  In its place they prove a different
characterization: in the general learning setting, a problem is learnable
if and only if it admits a \emph{stable}, asymptotically
empirical-risk-minimizing learner.  Uniform convergence is demoted from
characterization to a sufficient condition.

\medskip\noindent
\emph{What the witness exploits.}
The symmetrization argument at the heart of the Fundamental Theorem's
proof requires the loss to take finitely many values.  Specifically:
two.  That assumption is invisible in the binary setting: it looks like
a harmless feature of the problem.  Make the loss real-valued and
symmetrization breaks.  Uniform convergence stops characterizing
learnability; the stable-AERM condition takes over as the
characterization, indifferent to the cardinality of the loss.

The Fundamental Theorem is not wrong.  It is \emph{local}: an artifact
of the 0-1 loss that does not announce itself as such.

\hfill\cite{ShalevShwartzShamir2010}
\end{separation}


% ================================================================
% SECTION 3: THE STRICT STRENGTH HIERARCHY
% ================================================================

\section{Strict Strength Hierarchy}
\label{sec:strict-hierarchy}

The 4 \rel{strictly\_stronger} edges assert that one concept
\emph{does} imply another, but the reverse fails: the stronger concept
is a proper generalization. Each edge requires two proofs: one for the
forward implication, and one (the witness) for the strictness of the
gap.

% --- Strict 1 ---

\begin{separation}
\textbf{Ex-learning $\supsetneq$ FIN-learning.}\formal{FLT_Proofs.Theorem.Separation.ex_strictly_stronger_finite}

\medskip\noindent
\emph{Forward implication.} Every FIN-learnable class is trivially
Ex-learnable: a learner that outputs its final hypothesis after finitely
many data points and never changes it again is, a fortiori, a learner
that converges in the limit.

\medskip\noindent
\emph{Witness for strictness.}
FIN-learning requires zero mind changes: the learner must output its
final, correct hypothesis at some point and never retract it. Ex-learning
allows unbounded mind changes before convergence. The gap is witnessed
by any class requiring unbounded mind changes for identification.

Consider the class of pattern languages with growing structural
complexity. A Gold-style learner can identify any member in the limit
(Ex-learn it) by successively refining its conjecture as new data
arrives, but no learner can commit to a correct hypothesis after seeing
only finitely many initial elements without ever revising. The revision
process is essential, and FIN-learning forbids it.

In the Gold identification hierarchy, this separation is the first step:
$\FIN \subsetneq \Ex$, and the gap is measured by the mind-change
ordinal. FIN corresponds to $0$ mind changes; Ex allows $< \omega$ mind
changes (any finite number, not bounded in advance).

\hfill\cite{Gold1967}
\end{separation}

% --- Strict 2 ---

\begin{separation}
\textbf{Online learning $\supsetneq$ PAC learning.}\formal{FLT_Proofs.Theorem.Separation.online_strictly_stronger_pac}

\medskip\noindent
\emph{Forward implication.}
If $\mathcal{H}$ has finite Littlestone dimension $\Ldim(\mathcal{H})
= d$, then in particular $\VC(\mathcal{H}) \leq d$.\formal{FLT_Proofs.Theorem.Separation.vcdim_le_of_mistake_bounded} (Every shattered
set in the VC sense extends to a path in a Littlestone tree, so a
mistake-bounded learner with bound $d$ forces $\VC \leq d$.) By the
fundamental theorem, $\mathcal{H}$ is PAC learnable.

\medskip\noindent
\emph{Witness for strictness.}
Thresholds on finite domains provide the simplest witness. Fix
$X = \{1, 2, \ldots, n\}$ and let $\mathcal{H}_n = \{h_\theta :
\theta \in \{0, 1, \ldots, n\}\}$ where $h_\theta(x) = \ind[x >
\theta]$. Then $\VC(\mathcal{H}_n) = 1$ for every $n$, but
$\Ldim(\mathcal{H}_n) = \lfloor \log_2(n+1) \rfloor$, which grows
without bound as $n \to \infty$.

More dramatically, thresholds on $\R$ (the witness from Separation~1)
have $\VC = 1$ but $\Ldim = \infty$: PAC learnable but not online
learnable with any finite mistake bound. The containment is strict at
every level of the hierarchy.

\hfill\cite{Littlestone1988}
\end{separation}

% --- Strict 3 ---

\begin{separation}
\textbf{Universal learning $\supsetneq$ PAC learning.}

\medskip\noindent
\emph{Forward implication.}
Universal learnability is the \emph{weaker} requirement: it asks only for a
learner whose error decays at some rate under each distribution
\emph{separately}, whereas PAC learning demands a single sample-complexity
bound uniform over all distributions. Removing the uniformity, every
PAC-learnable class (finite VC dimension) is therefore universally
learnable, indeed at an exponential rate. The strictness below shows the
converse fails.

\medskip\noindent
\emph{Witness for strictness.}
Bousquet et al.~\cite{BousquetHannekeMoranVanHandelYehudayoff2021}
established a trichotomy for universal learning rates: every concept
class has either an exponential rate, an arbitrarily slow rate, or is
not universally learnable at all. The trichotomy theorem shows that
classes with $\VC(\mathcal{H}) = \infty$ but containing no infinite
Littlestone tree achieve exponential universal learning rates; they
are universally learnable but \emph{not} PAC learnable (since PAC
requires finite VC dimension).

The critical structural difference is the quantifier order.
PAC learning demands: $\exists$ learner $\forall$ distributions
$\forall \varepsilon, \delta$, the learner succeeds. Universal learning
demands: $\forall$ distributions $\exists$ rate at which \emph{some}
learner succeeds. By removing the uniformity requirement over
distributions, universal learning captures a strictly larger class
of learning problems.

\hfill\cite{BousquetHannekeMoranVanHandelYehudayoff2021}
\end{separation}

% --- Strict 4 ---

\begin{separation}
\textbf{DS dimension $\supsetneq$ Natarajan dimension.}

\medskip\noindent
\emph{Forward implication.}
The DS dimension refines the Natarajan dimension: $\Ndim(\mathcal{H})
\leq \DSdim(\mathcal{H})$ for every hypothesis class $\mathcal{H}$.
Both measure the combinatorial complexity of multiclass concept classes,
but the DS dimension accounts for the full oriented hypergraph structure,
not just pairwise shattering.

\medskip\noindent
\emph{Witness for strictness.}
Brukhim et al.~\cite{BrukhimCarmonDinurMoranYehudayoff2022} constructed
a concept class using a hyperbolic pseudo-manifold with $\Ndim = 1$
but $\DSdim = \infty$. The construction builds a concept class over an
infinite label space where every pair of points admits only one
non-trivial two-coloring (giving $\Ndim = 1$), but the global
combinatorial structure (encoded in the oriented faces of the
pseudo-manifold) is rich enough to drive the DS dimension to infinity.

This is the witness that simultaneously demonstrates the separation
$\Ndim \not\Rightarrow \text{PAC learnable}$ from \Cref{sec:separations}:
the same class has $\Ndim = 1$ but is not learnable, because learnability
is characterized by the DS dimension, not the Natarajan dimension.

\hfill\cite{BrukhimCarmonDinurMoranYehudayoff2022}
\end{separation}


% ================================================================
% SECTION 4: WHAT THE NEGATIVE LAYER REVEALS
% ================================================================

\section{What the Negative Layer Reveals}
\label{sec:neg-layer-summary}

The separation lattice of \Cref{fig:separation-lattice} carries a
meta-theorem about the structure of formal learning theory, which we
can now state explicitly.

\begin{remark}[The separation lattice is sparse]
\label{thm:sparsity}
This chapter establishes thirteen edges with explicit witnesses among the
major learning paradigms and complexity measures of
\Cref{fig:separation-lattice}, far fewer than the $\binom{k}{2}$ relations
that $k$ such concepts could \emph{a priori} exhibit.  We do not claim a
complete census of the lattice; the point is that the witnessed edges are
scarce, and that many remaining pairs appear unrelated: they operate on
different mathematical types, or are connected only by non-implicational
relations (analogy, measurement, assumption) rather than by an implication
that could be witnessed or refuted.
\end{remark}

This scarcity is not merely an artifact of incomplete knowledge.  Many of
the major paradigms operate on different mathematical types: PAC learning
and Gold-style identification take different data models (i.i.d.\ samples
versus infinite enumerations); online and exact learning use different
interaction protocols (adversarial instances versus query access); VC and
SQ dimension measure different properties (combinatorial shattering versus
correlation structure).  Between such concepts a clean implication is often
not even the right question.  Where we \emph{can} ask it, this chapter
witnesses one direction (Gold identification does not imply PAC learning,
say), while establishing the reverse non-implication is, in several cases, a
separate open problem rather than a settled fact.

The witnesses in this chapter make the incomparability concrete.
Thresholds on $\R$ separate PAC from online learning. All finite subsets
of $\N$ separate Gold identification from PAC learning. Parities
separate VC dimension from SQ dimension. Each witness exploits a specific
structural mismatch in the data model, the success criterion, the
adversarial model, or the computational model; and these mismatches are
not removable by clever proof techniques. They are features of the
mathematical landscape.

The four strict strength edges, taken together, form two chains:
\[
\text{Universal} \supsetneq \text{Online} \supsetneq \text{PAC}
\qquad \text{and} \qquad
\DSdim \supsetneq \Ndim.
\]
The first chain orders three paradigms by the stringency of their
learnability requirements. The second orders two dimensions by their
sensitivity to multiclass structure. In both cases, the strict
containment is proved by a single, explicit construction.

These are the theorems that textbooks usually omit. They deserve
their chapter.


\subsection{Where Attention Sits}
\label{sec:attention-lattice}

A separation lattice is a map, and it is natural to ask where a modern
hypothesis class lands on it. The softmax-attention concept class of
\Cref{ch:generalization} is machine-checked to be
well-behaved,\formal{TLT.Bridge.softAttention_wellBehaved} and with
norm-bounded parameters it has finite capacity
(\Cref{prop:certified-attention}); it therefore sits on the PAC node: it
is PAC learnable.

Its position relative to the \emph{online} node is open.  The
witness for Separation~1 was a class of finite VC dimension and infinite
Littlestone dimension (thresholds on~$\R$).  Attention compares query--key
scores, and threshold behavior on a score is exactly the structure that
sends the Littlestone dimension to infinity while the VC dimension stays
finite.  It is therefore possible that the softmax-attention class is its
\emph{own} witness for $\text{PAC} \not\Rightarrow \text{online}$, joining
thresholds on that dashed red edge.  Whether it is so is
\Cref{op:attention-littlestone}; the machine-checked
separation~\leanname{pac_not_implies_online} already supplies the template
a proof would instantiate.


% ================================================================
% OPEN PROBLEMS
% ================================================================

\section{Open Problems}
\label{sec:separations-open}

A separation chapter generates two kinds of open problem: the missing
\emph{reverse} of a one-sided non-implication, and the missing
\emph{formal proof} of a separation so far only cited.  Each problem below
is tagged a \emph{known unknown} (KU), sharply posed but unresolved, or an
\emph{unknown known} (UK), where the right object or invariant is still
missing; a design-lab or literature anchor is named where one exists.

\begin{openprob}[The Littlestone dimension of attention]
\label{op:attention-littlestone}
Is the Littlestone dimension of the softmax-attention concept class finite
or infinite?  If infinite, attention is its own witness for
$\text{PAC} \not\Rightarrow \text{online}$
(\leanname{pac_not_implies_online}); if finite, attention is online
learnable and the score comparisons are tamer than thresholds.  A known
unknown on the transformer side (\leanname{softAttention_wellBehaved}).
\end{openprob}

\begin{openprob}[Does PAC imply Ex?]
\label{op:reverse-pac-ex}
The chapter proves Gold identification does not imply PAC learning
(\leanname{ex_not_implies_pac}), but never the reverse.  Is there a
PAC-learnable concept class that is not Ex-learnable as a class of
functions, so that PAC and Gold are genuinely incomparable rather than
one-sidedly separated?  A known unknown.
\end{openprob}

\begin{openprob}[Formalizing the DS/Natarajan gap]
\label{op:formalize-ds-natarajan}
The pseudo-manifold witness with $\Ndim = 1$ and $\DSdim = \infty$ is cited
from Brukhim et al., not derived.  Can the construction, and the
non-learnability it implies, be machine-checked?  An unknown known: the
oriented-hypergraph combinatorics has no formal
home~\cite{BrukhimCarmonDinurMoranYehudayoff2022}.
\end{openprob}

\begin{openprob}[Formalizing the universal-learning trichotomy]
\label{op:formalize-trichotomy}
The forward edge universal$\,\Rightarrow\,$PAC is machine-checked, but the
trichotomy that witnesses Strict~3 (exponential, arbitrarily slow, or not
at all) is not.  Can the Bousquet et al.\ trichotomy be formalized,
closing the strictness direction?  An unknown
known~\cite{BousquetHannekeMoranVanHandelYehudayoff2021}.
\end{openprob}

\begin{openprob}[VC versus SQ, quantitatively]
\label{op:sq-vc-quantitative}
Parities give $\VC = n$, $\SQdim = 2^n$ under the uniform distribution.
Across all distributions, what is the tightest possible relationship
between VC and SQ dimension: is the exponential gap always achievable, or
distribution-dependent?  A known
unknown~\cite{BlumFurstJacksonKearnsMansourRudich1994}.
\end{openprob}

\begin{openprob}[Formalizing labeled versus unlabeled compression]
\label{op:formalize-compression-sep}
The P\'alv\"olgyi--Tardos $\VC = 2$ class separating labeled from unlabeled
size-2 compression is cited, not constructed here.  Can the separation be
machine-checked, and does it persist at larger VC dimension and scheme
size?  An unknown known~\cite{PalvolgyiTardos2020}.
\end{openprob}

\begin{openprob}[Is the lattice complete?]
\label{op:lattice-completeness}
For each pair of paradigms with no witnessed edge, is the pair provably
incomparable, or merely unstudied?  A census distinguishing genuine
incomparability from undiscovered implications would turn the informal
sparsity of \Cref{thm:sparsity} into a theorem.  A known unknown.
\end{openprob}

\begin{openprob}[Characterizing the proper/improper gap]
\label{op:proper-improper-sep}
3-DNF has an efficient improper learner but no efficient proper one
(\leanname{IsProper}, under $\mathrm{RP}\neq\mathrm{NP}$).  Which structural
property of a class forces a proper/improper efficiency gap, and which
classes have none?  A known unknown~\cite{PittValiant1988}.
\end{openprob}

\begin{openprob}[The scope of the stability characterization]
\label{op:stability-scope}
Stable AERM characterizes learnability in the general learning setting
(\Cref{sec:separations}, Separation~9).  Exactly which loss families and
hypothesis structures fall inside that characterization, and where does
even stability cease to characterize?  A known
unknown~\cite{ShalevShwartzShamir2010}.
\end{openprob}

\begin{openprob}[The SQ dimension of attention]
\label{op:attention-sq}
Parities resist statistical-query learning despite small VC.  Does the
softmax-attention class behave like a VC class under SQ, or can it encode
parity-like correlations that make it SQ-hard?  An unknown known on the
transformer side (\leanname{softAttention_wellBehaved}).
\end{openprob}
